\section{What complex analysis means in iPRA}\label{sec:ipra}

\subsection{The object theory}

The repository snapshot used here is commit
\code{85afffc93b2abbc1972fdfc98244052cbfbda1b3}.  Its ordinary presentation of
iPRA has first-order intuitionistic logic with equality, a symbol for every
primitive-recursive function, the defining equations for those symbols,
successor axioms, and induction for quantifier-free formulas.  In Lean these
are the types \code{PR}, \code{PTm}, and \code{PRAAx}; the induction constructor
is \code{PRA.qfInd}.  Consequently, a proof may use an externally fixed
primitive-recursive approximation algorithm, but it may not quantify over all
algorithms by invoking a universal evaluator.

This distinction controls the formulation of every theorem below.  A statement
such as ``for every complex parameter $a$'' is read as a schema: fix a
primitive-recursive name for $a$ and the certificates demanded by the theorem;
then iPRA proves the resulting first-order sentence.  For rational parameters
the package also gives a single uniform coded algorithm.

\subsection{Rational and complex codes}

A rational is represented by a sign and two natural numbers $(u,v)$ with
$v>0$; normalization is optional.  Equality is the cross-multiplied equation
$u_1v_2=u_2v_1$, so no gcd computation is logically required.  A rational
complex number is a pair of rational codes.  We use the $\ell^1$ magnitude
\[
 \norm{x+\ii y}_1=\abs{x}+\abs{y},
\]
because all its upper-bound relations reduce to integer inequalities.  It is
submultiplicative up to the exact estimate
$\norm{zw}_1\le\norm z_1\norm w_1$.

\begin{definition}[Dyadic ball equality]
For rational complex codes $u,v$ and $k\in\N$,
\[
 \BallEq_k(u,v)\quad:\Longleftrightarrow\quad
 \norm{u-v}_1\le 2^{-k}.
\]
After clearing positive denominators, this is quantifier-free arithmetic.
\end{definition}

A complex name is an externally fixed primitive-recursive function
$u:\N\to\Q(\ii)$ satisfying
\begin{equation}\label{eq:name-coherence}
 \BallEq_k(u(k+2),u(k+1))
 \quad\text{and}\quad
 \BallEq_k(u(k+1),u(k)).
\end{equation}
After a harmless reindexing, the geometric series proves that all sufficiently
accurate approximants represent one observational value.  No completed Cauchy
real sort is added to the language.

\subsection{Apartness, branches, and safe division}

Division by $w$ is permitted only with an apartness witness
$\norm w_1\ge2^{-r}$.  A branch certificate for $z$ consists of:
\begin{enumerate}
\item an apartness witness $\norm z_1\ge2^{-r}$;
\item a rational angular interval $I$ of width below $2\pi$ containing the
      argument of every allowed approximation to $z$;
\item a rational clearance from the chosen branch cut;
\item primitive-recursive approximants to $\Log z$ with a modulus derived from
      those clearances.
\end{enumerate}
Then $z^a$ means $\exp(a\Log z)$ on that named branch.  All identities involving
powers use the same branch on both sides.  This removes a frequent hidden
classical error: $z^{a+b}=z^az^b$ is not a branch-free identity.

For the positive real special cases, a branch certificate is just a positive
rational lower bound, and the construction simplifies substantially.

\subsection{Finite analytic contracts}

\begin{contract}[Gamma]\label{con:gamma}
For a fixed compact rational box $K$ with rational distance $\eta>0$ from the
nonpositive integers, \code{GammaApprox} returns, from $(w,k)$, a rational
complex ball of radius $2^{-k}$ containing $\Gamma(w)$ uniformly for $w\in K$.
It also returns derivative jets of any externally fixed order and vertical-line
bounds
\begin{align}
 \abs{\Gamma(\sigma+\ii t)}
 &\le C(1+\abs t)^{\sigma-1/2}\e^{-\pi\abs t/2},\label{eq:gamma-upper}\\
 \abs{1/\Gamma(\sigma+\ii t)}
 &\le C'(1+\abs t)^{1/2-\sigma}\e^{\pi\abs t/2},\label{eq:recip-gamma-upper}
\end{align}
for $\abs t\ge T_0$, with rational $C,C',T_0$ computed from $K$ and $\eta$.
\end{contract}
The existing \code{gamma\_ipra} packet develops the finite Euler products and
Stirling bounds needed for this contract.  That packet itself explicitly does
not claim a completed analytic object derivation; our formalization plan treats
\cref{con:gamma} as a module interface to be discharged.

\begin{contract}[Power series]
A series $\sum c_r$ is certified by a primitive-recursive index $N(k)$ and an
upper bound $E(N)$ such that finite arithmetic proves
$\sum_{r\ge N}\norm{c_r}_1\le E(N)\le2^{-k}$.  Differentiated series carry the
same data for each externally fixed derivative order.
\end{contract}

\begin{contract}[Contour integral]
A contour is a rational polygonal path, or a vertical segment with rational
real part, together with: a lower distance from every pole and branch cut; a
bound on the integrand and a fixed number of derivatives on each segment; a
finite quadrature rule; and endpoint/tail errors.  The value at precision $k$
is a finite rational sum with total error at most $2^{-k}$.
\end{contract}

\begin{contract}[Improper real integral and Fubini]
An integral on $(0,\infty)$ is split at rational $0<\epsilon<1<R$.  The middle
piece has a finite quadrature certificate; the two endpoint pieces have
explicit majorants.  A change of integration order additionally supplies a
majorant for the discarded rectangles and an error budget proving that both
orders approximate the same finite double sum.
\end{contract}

All four contracts are arithmetical data.  For fixed analytic names, their
verification consists of bounded loops and quantifier-free induction.

\subsection{Formula shape in the object language}

A typical analytic theorem is not written inside iPRA as an equality of
abstract complex numbers.  Its precision-$k$ instance has the form
\begin{equation}\label{eq:object-shape}
 \forall x\le B(k)\;
   \bigl(\Adm(d,x,k)\to
   \BallEq_k(\Eval_{\rm left}(d,x,k+c),
             \Eval_{\rm right}(d,x,k+c))\bigr),
\end{equation}
where $c$ is a fixed guard-bit budget.  Every symbol in
\eqref{eq:object-shape} is primitive recursive.  Unbounded parameters are
handled by a schema in the external choices $d$ and $B$, not by a universal
program interpreter.

The finite triangle inequality gives transitivity of ball equality.  Thus a
proof graph may pass from a contour evaluator to a residue evaluator, then to a
hypergeometric evaluator, and finally to an elementary evaluator.  The total
precision budget is allocated backwards through that graph.  This is the
mechanism by which the many special cases become ordinary iPRA derivations.
